\documentclass[11pt]{article}

% ---------------------------------------------------------------------------
% Working paper: mathematical companion to the jump-diffusion-estimation
% Python library. Pedagogical in the spirit of Ospina Arango (2009): every
% method is derived mathematically and then grounded in the actual code
% that implements it.
% ---------------------------------------------------------------------------

\usepackage[T1]{fontenc}
\usepackage[utf8]{inputenc}
\usepackage{lmodern}
\usepackage{microtype}
\usepackage{amsmath,amssymb,amsthm}
\usepackage[margin=2.8cm]{geometry}
\usepackage{booktabs}
\usepackage{xcolor}
\usepackage{listings}
\usepackage[round]{natbib}
\usepackage[colorlinks=true,linkcolor=blue!50!black,citecolor=blue!50!black,urlcolor=blue!50!black]{hyperref}

% ----- Code listings ---------------------------------------------------------
\definecolor{codegreen}{rgb}{0.0,0.45,0.0}
\definecolor{codegray}{rgb}{0.5,0.5,0.5}
\definecolor{codeblue}{rgb}{0.15,0.15,0.6}
\definecolor{codebg}{rgb}{0.97,0.97,0.97}

\lstdefinestyle{pystyle}{
  language=Python,
  backgroundcolor=\color{codebg},
  commentstyle=\color{codegreen}\itshape,
  keywordstyle=\color{codeblue}\bfseries,
  stringstyle=\color{codegreen},
  numberstyle=\tiny\color{codegray},
  basicstyle=\ttfamily\small,
  breaklines=true,
  breakatwhitespace=true,
  showstringspaces=false,
  numbers=left,
  numbersep=8pt,
  frame=single,
  framerule=0.3pt,
  rulecolor=\color{codegray},
  xleftmargin=16pt,
  aboveskip=10pt,
  belowskip=10pt,
  captionpos=b
}
\lstset{style=pystyle}

% ----- Theorem-like environments ---------------------------------------------
\emergencystretch=1.5em

\theoremstyle{definition}
\newtheorem{definition}{Definition}[section]
\newtheorem{example}{Example}[section]
\theoremstyle{remark}
\newtheorem{remark}{Remark}[section]
\newtheorem{decision}{Design decision}[section]

% ----- Notation shortcuts -----------------------------------------------------
\newcommand{\dd}{\,\mathrm{d}}
\newcommand{\E}{\mathbb{E}}
\newcommand{\Prob}{\mathbb{P}}
\newcommand{\R}{\mathbb{R}}
\newcommand{\Ind}[1]{\mathbf{1}\{#1\}}
\newcommand{\code}[1]{\texttt{#1}}
\newcommand{\loglik}{\ell}
\newcommand{\thetavec}{\vartheta}

\title{%
  Likelihood-Based Estimation of Jump-Diffusion Processes in Python:\\
  The Mathematics Behind the \code{jump-diffusion-estimation} Library\\[4pt]
  {\large A pedagogical working paper}%
}
\author{%
  Juan David Ospina Arango\thanks{%
    Departamento de Ciencias de la Computaci\'on y de la Decisi\'on,
    Facultad de Minas, Universidad Nacional de Colombia, Medell\'in, Colombia.
    \code{jdospina@gmail.com}. ORCID: 0000-0002-3547-5247.}%
}
\date{Working draft --- \today}

\begin{document}

\maketitle

\begin{abstract}
\noindent
This working paper documents, in a deliberately pedagogical style, the
mathematics implemented by the open-source Python library
\code{jump-diffusion-estimation}. The library simulates and estimates
discretely sampled jump-diffusion processes whose jump sizes follow a
pluggable distribution (normal, skew-normal, skewed generalized error,
Kou double-exponential, Student-$t$). We derive each component in the order
a user encounters it: the stochastic differential equation and its
Bernoulli (at-most-one-jump) transition density; the jump-size families and
the closed-form Gaussian convolutions available for some of them; the
FFT-based numerical convolution used when no closed form exists; maximum
likelihood estimation with both a local optimizer (L-BFGS-B) and a global
one (differential evolution); three uncertainty-quantification routes
(profile likelihood, Wald via the observed Fisher information, and the
parametric bootstrap); a parametric-bootstrap likelihood-ratio test for the
presence of jumps; and goodness-of-fit comparison across jump
distributions. Every derivation is followed by the actual code that
implements it, together with the design decisions --- grid sizes,
penalization devices, optimizer defaults --- and the reasons behind them.
The paper is a companion to the library and to the author's 2009 thesis,
from which several of the numerical schemes are ported.
\end{abstract}

\vspace{1em}
\noindent\textbf{Keywords:} jump-diffusion; maximum likelihood; FFT
convolution; differential evolution; profile likelihood; parametric
bootstrap; Python.

\tableofcontents

\section{Introduction}
\label{sec:introduction}

Financial asset returns exhibit occasional movements far too large to be
explained by a Gaussian diffusion. \citet{Merton1976} formalized the now
classical remedy: superimpose a compound jump process on the Brownian
diffusion, so that ordinary days are driven by the diffusion and
extraordinary days by the jumps. Fitting such a model to discretely
sampled data by maximum likelihood is conceptually straightforward ---
write the transition density of an increment, take logs, sum, maximize ---
but almost every step hides a numerical or statistical subtlety:

\begin{itemize}
  \item the transition density involves the \emph{convolution} of the
        Gaussian diffusion with the jump-size law, which has a closed form
        only for a few jump distributions;
  \item the resulting mixture likelihood is \emph{multimodal}, so standard
        gradient-based optimizers can converge to spurious local optima or
        stall at box constraints;
  \item the numerical Hessian at the optimum is frequently ill-conditioned,
        which makes the textbook (Wald) standard errors fragile;
  \item testing whether jumps are present at all puts the null hypothesis
        on the \emph{boundary} of the parameter space while leaving the
        jump-size parameters \emph{unidentified}, so the classical
        likelihood-ratio theory does not apply.
\end{itemize}

The Python library
\code{jump-diffusion-estimation}\footnote{\url{https://github.com/jdospina/jump-diffusion-estimation}}
packages a coherent solution to each of these problems. This working paper
explains \emph{how}: every section states the mathematics first and then
shows the code that implements it, including the design decisions ---
discretization grids, penalization devices, optimizer configurations ---
and the reasoning behind them. Several of those decisions are direct ports
of the author's MSc thesis \citep{Ospina2009} and of the subsequent applied
work in \citet{Ardia2011}; where that is the case we say so explicitly, so
the reader can trace each formula to its origin.

\subsection{The library at a glance}
\label{sec:architecture}

The package is organized in four cooperating layers, each mapped to a
Python subpackage:

\begin{center}
\small
\begin{tabular}{@{}llp{5.4cm}@{}}
\toprule
Layer & Subpackage & Main classes \\
\midrule
Model \& simulation & \code{models}, \code{simulation} &
  \code{JumpDiffusionModel},\newline \code{JumpDiffusionSimulator} \\
Jump distributions & \code{distributions} &
  \code{JumpDistribution} + five concrete families \\
Estimation \& inference & \code{estimation} &
  \code{JumpDiffusionEstimator} \\
Validation \& comparison & \code{validation} &
  \code{ValidationExperiment},\newline \code{JumpDistributionComparison} \\
\bottomrule
\end{tabular}
\end{center}

The dependency direction is strictly downward in the table:
\code{JumpDiffusionModel} owns the transition density and path simulation,
delegating everything jump-specific to a \code{JumpDistribution} object;
\code{JumpDiffusionEstimator} evaluates that model's likelihood at trial
parameter vectors and wraps the optimizers and the inference machinery;
the validation layer repeatedly drives the estimator on simulated or real
data. A user who wants a new jump-size family implements a single class
with the distribution's density --- everything else (likelihood, FFT
convolution, sampling, estimation, inference, comparison) works unchanged.
That extension contract is the subject of Section~\ref{sec:distributions}.

\subsection{Notation}
\label{sec:notation}

Throughout the paper, $X_t$ denotes the (log-price) process, $\mu$ the
drift, $\sigma > 0$ the diffusion volatility, $W_t$ a standard Brownian
motion, $N_t$ a Poisson counting process with intensity $\lambda$, and
$J$ a generic jump size with density $f_J(\cdot\,;\theta_J)$ and parameter
vector $\theta_J$. Observations are $n$ increments
$\Delta x_i = x_{t_i} - x_{t_{i-1}}$ on a regular grid with spacing
$\Delta t$; in the code, \code{data} is the array of increments and
\code{dt} is $\Delta t$. The Bernoulli jump probability per interval is
$p$ (\code{jump\_prob} in the code), which plays the role of
$\lambda \Delta t$. The full parameter vector is
\[
  \thetavec \;=\; (\mu,\ \sigma,\ p,\ \theta_J),
\]
and the code orders it exactly this way:
\begin{center}
\code{("mu", "sigma", "jump\_prob", *jump\_distribution.param\_names)}.
\end{center}
We write $\phi(x; m, s^2)$ for the normal density with mean $m$ and
variance $s^2$, $\Phi$ for the standard normal CDF, $*$ for convolution,
and $\loglik(\thetavec)$ for the log-likelihood. Code identifiers are set
in \code{typewriter} font; all listings are excerpts of the actual library
source (docstrings and defensive branches are sometimes elided for
readability, never altered).

\subsection{How to read this paper}

Sections follow the pipeline of an actual analysis.
Section~\ref{sec:model} derives the model and its Bernoulli-approximated
transition density --- the single formula the whole library revolves
around. Section~\ref{sec:distributions} presents the five jump-size
families and their closed-form Gaussian convolutions where they exist.
Section~\ref{sec:fft} derives the FFT convolution scheme used when they
do not. Section~\ref{sec:estimation} covers maximum likelihood and the
two optimizers. Section~\ref{sec:inference} derives the three
uncertainty-quantification routes, Section~\ref{sec:jumptest} the
bootstrap test for the presence of jumps, and Section~\ref{sec:comparison}
the goodness-of-fit comparison and Monte Carlo validation tools.
Section~\ref{sec:closing} closes with the research agenda this library
supports.

\section{The model and its transition density}
\label{sec:model}

\subsection{The stochastic differential equation}

The library implements the arithmetic jump-diffusion
\begin{equation}
  \dd X_t \;=\; \mu \dd t \;+\; \sigma \dd W_t \;+\; J_t \dd N_t,
  \label{eq:sde}
\end{equation}
where $N_t$ is a Poisson process with intensity $\lambda$, independent of
$W_t$, and the jump sizes $J_t$ are i.i.d.\ draws from
$f_J(\cdot\,;\theta_J)$, independent of both. Because the coefficients are
constant, \eqref{eq:sde} integrates in closed form:
\begin{equation}
  X_t \;=\; X_0 \;+\; \mu t \;+\; \sigma W_t \;+\; \sum_{k=1}^{N_t} J_k .
  \label{eq:solution}
\end{equation}
The arithmetic specification is intended for \emph{log-prices}: if $S_t$
is the price and $X_t = \log S_t$, then \eqref{eq:sde} is the standard
Merton model written on returns, and the increments
$\Delta x_i$ are the log-returns the user hands to the estimator.

\begin{remark}
Equation \eqref{eq:solution} is exact, not an Euler approximation ---
constant coefficients make the integral trivial. Every approximation in
the library enters later, through the treatment of the jump count per
observation interval, not through time discretization of the diffusion.
\end{remark}

\subsection{Exact transition density: the Poisson mixture}

Over an interval of length $\Delta t$, the increment implied by
\eqref{eq:solution} is
\begin{equation}
  \Delta X \;=\; \mu\,\Delta t \;+\; \sigma \sqrt{\Delta t}\, Z
  \;+\; \sum_{k=1}^{K} J_k,
  \qquad Z \sim \mathcal N(0,1),\quad
  K \sim \mathrm{Poisson}(\lambda \Delta t),
  \label{eq:increment}
\end{equation}
with $Z$, $K$ and the $J_k$ mutually independent. Conditioning on
$K = k$ and convolving gives the exact transition density as a Poisson
mixture:
\begin{equation}
  g(\Delta x)
  \;=\;
  \sum_{k=0}^{\infty}
  e^{-\lambda \Delta t}\,\frac{(\lambda \Delta t)^k}{k!}
  \left( \phi_{\mu \Delta t,\, \sigma^2 \Delta t} * f_J^{*k} \right)
  (\Delta x),
  \label{eq:poisson-mixture}
\end{equation}
where $f_J^{*k}$ is the $k$-fold convolution of the jump density
($f_J^{*0} = \delta_0$). Each additional term requires one more
convolution, and for most jump laws none of them is available in closed
form.

\subsection{The Bernoulli approximation}
\label{sec:bernoulli}

\begin{decision}[At most one jump per interval]
\label{dec:bernoulli}
For daily data, $\lambda \Delta t$ is small: at $\lambda \Delta t = p$,
the probability of two or more jumps in one interval is
\[
  \Prob(K \ge 2) \;=\; 1 - e^{-p}(1 + p) \;=\; \tfrac{p^2}{2} + O(p^3),
\]
i.e.\ about $0.005$ even at the fairly extreme $p = 0.1$. The library
therefore truncates \eqref{eq:poisson-mixture} at $k \le 1$ and replaces
the Poisson count by a Bernoulli indicator $B \sim \mathrm{Bernoulli}(p)$
with $p = \lambda \Delta t$:
\begin{equation}
  \Delta X \;\approx\; \mu\,\Delta t + \sigma \sqrt{\Delta t}\, Z + B\,J .
  \label{eq:bernoulli-increment}
\end{equation}
This is the same approximation used in \citet{Ospina2009} and
\citet{Ardia2011}. It buys a two-component mixture density (one
convolution instead of infinitely many) at a bias cost of order
$(\lambda\Delta t)^2$; quantifying that bias precisely across the
$\lambda \Delta t$ range is part of the research agenda of
Section~\ref{sec:closing}.
\end{decision}

Under \eqref{eq:bernoulli-increment}, the transition density becomes the
two-component mixture at the heart of the library,
\begin{equation}
  \boxed{\;
  g_{\thetavec}(\Delta x)
  \;=\;
  (1 - p)\,\phi\!\left(\Delta x;\ \mu\Delta t,\ \sigma^2 \Delta t\right)
  \;+\;
  p \left( \phi_{\mu\Delta t,\,\sigma^2\Delta t} * f_J \right)(\Delta x).
  \;}
  \label{eq:mixture}
\end{equation}
Interpreting \eqref{eq:mixture} reads exactly like the data-generating
story: with probability $1-p$ the day is ``ordinary'' and the increment is
Gaussian; with probability $p$ a jump occurred and the increment is the sum
of the Gaussian part and one draw from the jump law --- hence the
convolution.

\subsection{The code: simulation}

\code{JumpDiffusionModel.simulate} generates paths directly from
\eqref{eq:bernoulli-increment}: a vector of standard normal innovations, a
vector of Bernoulli jump indicators, and a vector of jump sizes drawn from
the jump distribution, combined step by step.

\begin{lstlisting}[caption={Path simulation in \code{models/jump\_diffusion.py} (excerpt).}]
dt = T / n_steps
brownian_innovations = np.random.normal(0, 1, n_steps)
jump_indicators = np.random.binomial(1, self.parameters["jump_prob"], n_steps)
jump_sizes = self.jump_distribution.rvs(self._jump_params(), size=n_steps)

for i in range(n_steps):
    drift = self.parameters["mu"] * dt
    sigma_term = self.parameters["sigma"] * np.sqrt(dt)
    diffusion = sigma_term * brownian_innovations[i]
    jump = jump_indicators[i] * jump_sizes[i]
    path[i + 1] = path[i] + drift + diffusion + jump
\end{lstlisting}

Note that a jump size is drawn for every step but only \emph{applied} when
the indicator fires; this keeps the random-number stream aligned across
parameter values, which is convenient for variance-reduction and
reproducibility in Monte Carlo work.

\subsection{The code: the mixture density}

The model's \code{\_mixture\_density} method evaluates \eqref{eq:mixture}.
The only branch is \emph{how} the convolution term is computed: the jump
distribution's \code{diffusion\_convolved\_pdf} is first asked for a
closed form (Section~\ref{sec:distributions}); if it returns \code{None},
the generic FFT fallback of Section~\ref{sec:fft} takes over.

\begin{lstlisting}[caption={The mixture density \eqref{eq:mixture} in \code{models/jump\_diffusion.py}.}]
def _mixture_density(self, x: np.ndarray, dt: float) -> np.ndarray:
    mu = self.parameters["mu"]
    sigma = self.parameters["sigma"]
    p = self.parameters["jump_prob"]

    diffusion_mean = mu * dt
    diffusion_std = sigma * np.sqrt(dt)
    diffusion_density = norm.pdf(x, loc=diffusion_mean, scale=diffusion_std)

    # Jump + diffusion component: use a closed form if the jump
    # distribution has one, otherwise fall back to FFT convolution.
    jump_params = self._jump_params()
    jump_diffusion_density = self.jump_distribution.diffusion_convolved_pdf(
        x, jump_params, diffusion_mean, diffusion_std
    )
    if jump_diffusion_density is None:
        jump_diffusion_density = self.jump_distribution.fft_convolved_pdf(
            x, jump_params, diffusion_mean, diffusion_std
        )

    return (1 - p) * diffusion_density + p * jump_diffusion_density
\end{lstlisting}

The log-likelihood of the observed increments is then simply
\begin{equation}
  \loglik(\thetavec)
  \;=\;
  \sum_{i=1}^{n} \log g_{\thetavec}(\Delta x_i),
  \label{eq:loglik}
\end{equation}
computed in the model's \code{log\_likelihood} method with one numerical
guard: densities are floored at $10^{-300}$ before taking logs, so a
single increment falling in a region where the numerical density
underflows to zero degrades the likelihood instead of destroying it with
$\log 0 = -\infty$.

\begin{lstlisting}[caption={Log-likelihood \eqref{eq:loglik} with underflow guard.}]
def log_likelihood(self, data: np.ndarray, dt: float) -> float:
    increments = data if data.ndim == 1 else np.diff(data)
    densities = self._mixture_density(increments, dt)
    densities = np.maximum(densities, 1e-300)  # Numerical stability
    return np.sum(np.log(densities))
\end{lstlisting}

\section{Jump-size distributions}
\label{sec:distributions}

Everything jump-specific in the library lives behind one abstract
interface, \code{JumpDistribution}. This section first describes that
contract, then walks through the five concrete families, deriving the
closed-form Gaussian convolution wherever one exists --- those closed
forms are what let the likelihood skip the FFT machinery of
Section~\ref{sec:fft} for the families that admit them.

\subsection{The extension contract}
\label{sec:contract}

A concrete jump distribution must implement four small methods and may
override three more:

\begin{center}
\begin{tabular}{@{}llp{7.2cm}@{}}
\toprule
Method & Status & Role \\
\midrule
\code{pdf(x, params)} & required & jump-size density $f_J(x;\theta_J)$ \\
\code{default\_params()} & required & sensible defaults for simulation \\
\code{param\_bounds()} & required & optimization box for $\theta_J$ \\
\code{initial\_guess(m, s, skew)} & required & moment-based starting values \\
\midrule
\code{characteristic\_scale(params)} & optional & scale $s_J$ used to size numerical grids \\
\code{characteristic\_location(params)} & optional & center $\ell_J$ used to place numerical grids \\
\code{diffusion\_convolved\_pdf(...)} & optional & closed form for $\phi * f_J$, if known \\
\code{rvs(params, size)} & optional & native sampler, if available \\
\code{mean(params)}, \code{variance(params)} & generic & numerical $\E[J]$, $\operatorname{Var}[J]$ \\
\bottomrule
\end{tabular}
\end{center}

\begin{decision}[Density-only extensibility]
\label{dec:pdf-only}
The base class provides generic numerical fallbacks for the two hard
operations: \code{fft\_convolved\_pdf} approximates $\phi * f_J$ by FFT
(Section~\ref{sec:fft}), and \code{rvs} samples by numerically inverting
the CDF obtained from the cumulative sum of the discretized density. A new
distribution therefore only \emph{needs} its \code{pdf}; closed-form
overrides are pure speed-ups. This is what makes the library's ``pluggable
jump law'' promise cheap to keep.
\end{decision}

The numerical fallbacks share one discretization of the density,
\code{\_moment\_grid}: $2m = 2^{13}$ half-integer-offset points (coarser
than the FFT grid --- ample for sampling and moments) with step
$h = s_J / 200$, \emph{centered on the law's characteristic location}
$\ell_J$ so that a shifted distribution (a Student-$t$ or SGED with a
large \code{jump\_loc}) is not sampled from an empty grid. On top of it,
the generic \code{rvs} inverts the cumulative sum by linear
interpolation, and the generic \code{mean}/\code{variance} integrate the
discretized density --- the moments that feed the model-implied increment
diagnostics of Section~\ref{sec:inference}, verified in the test suite
against the closed-form moments of the normal, skew-normal, Kou and
located Student-$t$ laws:

\begin{lstlisting}[caption={Shared grid, inverse-CDF sampling and numerical moments in \code{distributions/base.py} (excerpts).}]
def _moment_grid(self, params):
    jump_std = self.characteristic_scale(params)
    h = jump_std / 200.0
    m = 2**12  # a coarser grid than fft_convolved_pdf suffices here
    k = np.arange(-m + 0.5, m, 1.0)
    # Centered on the distribution's location so that a shifted
    # distribution (large jump_loc) does not fall off the grid.
    x_grid = self.characteristic_location(params) + k * h
    density = np.maximum(self.pdf(x_grid, params), 0.0)
    return x_grid, density, h

def rvs(self, params, size, random_state=None):
    x_grid, density, h = self._moment_grid(params)
    cdf = np.cumsum(density) * h
    cdf /= cdf[-1]
    u = rand.uniform(0.0, 1.0, size=size)
    return np.interp(u, cdf, x_grid)

def mean(self, params):
    x_grid, density, h = self._moment_grid(params)
    total = float(np.sum(density) * h)
    if total <= 0:
        return float("nan")
    return float(np.sum(x_grid * density) * h / total)
\end{lstlisting}

The two \code{characteristic\_*} hooks deserve a word: the numerical
grids must be sized by how spread out the jump law is and placed where
its mass actually lives. The default \code{characteristic\_scale} reads
the parameter named \code{jump\_scale} (the convention of the normal,
skew-normal, SGED and Student-$t$ families); distributions parameterized
differently override it --- Kou's law, with separate up and down scales,
returns the larger of the two. The default
\code{characteristic\_location} reads \code{jump\_loc}, falling back to
$0$ for the families centered at the origin; laws whose center is not a
parameter but stays within a few scales of zero (Kou again) are covered
by the grids' tail margins and need no override.

\subsection{Normal jumps (Merton)}
\label{sec:normal}

The classical model of \citet{Merton1976}: $J \sim \mathcal N(0, s^2)$
with a single parameter \code{jump\_scale}~$= s$. Its convolution with the
diffusion is the textbook result that a sum of independent normals is
normal:
\begin{equation}
  \phi_{m_d,\,\sigma_d^2} * f_J
  \;=\;
  \phi\!\left(\cdot\,;\ m_d,\ \sigma_d^2 + s^2\right),
  \label{eq:normal-conv}
\end{equation}
where $m_d = \mu\Delta t$ and $\sigma_d = \sigma\sqrt{\Delta t}$ are the
diffusion mean and standard deviation of Section~\ref{sec:model}.

\begin{lstlisting}[caption={Closed-form convolution \eqref{eq:normal-conv} in \code{distributions/normal.py}.}]
def diffusion_convolved_pdf(self, x, params, diffusion_mean, diffusion_std):
    combined_std = np.sqrt(diffusion_std**2 + params["jump_scale"] ** 2)
    return norm.pdf(x, loc=diffusion_mean, scale=combined_std)
\end{lstlisting}

Beyond historical importance, \code{NormalJump} plays a structural role:
it is nested inside the skew-normal family (at zero skewness) and inside
the SGED (at $\nu = 2$, $\xi = 1$), which makes it the natural null model
when asking whether asymmetry or heavy tails are really needed
(Section~\ref{sec:comparison}).

\subsection{Skew-normal jumps}
\label{sec:skewnormal}

The skew-normal distribution of \citet{Azzalini1985} with location $0$,
scale $\omega$ (\code{jump\_scale}) and shape $\alpha$ (\code{jump\_skew})
has density
\begin{equation}
  f_J(x)
  \;=\;
  \frac{2}{\omega}\,
  \phi\!\left(\frac{x}{\omega}\right)
  \Phi\!\left(\alpha \frac{x}{\omega}\right).
  \label{eq:sn-pdf}
\end{equation}
Positive $\alpha$ skews the jumps to the right, negative to the left, and
$\alpha = 0$ recovers the normal.

The family is closed under convolution with an independent normal
\citep{Azzalini1986}\footnote{The library's code comments cite this
result as ``Azzalini \& Henze (1986)''; the closure property itself
appears in \citet{Azzalini1986}, while Henze's 1986 paper gives the
related probabilistic representation of the skew-normal.}: if $J \sim \mathrm{SN}(0, \omega, \alpha)$ and
$D \sim \mathcal N(m_d, \sigma_d^2)$ independently, then $D + J$ is again
skew-normal with location $m_d$, scale
$\tilde\omega = \sqrt{\omega^2 + \sigma_d^2}$, and a \emph{new} shape
parameter obtained by passing through the $\delta$-parameterization:
\begin{equation}
  \delta = \frac{\alpha}{\sqrt{1 + \alpha^2}},
  \qquad
  \tilde\delta = \frac{\omega\,\delta}{\tilde\omega},
  \qquad
  \tilde\alpha = \frac{\tilde\delta}{\sqrt{1 - \tilde\delta^{\,2}}}.
  \label{eq:sn-conv}
\end{equation}

\begin{decision}[Go through $\delta$, not $\alpha$]
\label{dec:delta}
It is tempting to guess that the new shape is simply
$\alpha\,\omega / \tilde\omega$ --- rescale the shape the way the scale was
rescaled. That shortcut is \emph{wrong}: it is off by about $0.2\%$ at
$\alpha = 2$, an error easily mistaken for FFT discretization noise but
which biases the likelihood systematically. The library's implementation
was verified against direct numerical convolution, and the code carries a
comment warning future maintainers away from the shortcut.
\end{decision}

\begin{lstlisting}[caption={Skew-normal convolution \eqref{eq:sn-conv} in \code{distributions/skew\_normal.py}.}]
def diffusion_convolved_pdf(self, x, params, diffusion_mean, diffusion_std):
    omega = params["jump_scale"]
    alpha = params["jump_skew"]
    combined_std = np.sqrt(diffusion_std**2 + omega**2)

    # Azzalini & Henze (1986): the new shape parameter is obtained by
    # going through the "delta" parameterization, not alpha*omega/std.
    delta = alpha / np.sqrt(1 + alpha**2)
    adjusted_delta = omega * delta / combined_std
    adjusted_skew = adjusted_delta / np.sqrt(1 - adjusted_delta**2)

    return skewnorm.pdf(x, a=adjusted_skew, loc=diffusion_mean,
                        scale=combined_std)
\end{lstlisting}

\subsection{SGED jumps}
\label{sec:sged}

The skewed generalized error distribution is the library's most flexible
symmetric-plus-skewness family and the centerpiece of
\citet{Ospina2009}; see also \citet{Theodossiou2015} for its role in
finance. It is built in three steps, and the code mirrors the construction
function by function.

\paragraph{Step 1: the GED base.}
The generalized error distribution with shape $\nu > 0$, standardized to
mean $0$ and variance $1$, has density
\begin{equation}
  f_{\mathrm{GED}}(z; \nu)
  \;=\;
  \frac{\nu \, \exp\!\left( -\tfrac12 \left| z / \lambda_\nu \right|^{\nu} \right)}
       {\lambda_\nu\, 2^{1 + 1/\nu}\, \Gamma(1/\nu)},
  \qquad
  \lambda_\nu = \sqrt{\frac{2^{-2/\nu}\,\Gamma(1/\nu)}{\Gamma(3/\nu)}} .
  \label{eq:ged}
\end{equation}
$\nu = 2$ gives the standard normal, $\nu = 1$ the Laplace; smaller $\nu$
means heavier tails.

\paragraph{Step 2: skewing \`a la Fern\'andez--Steel.}
\citet{FernandezSteel1998} skew any symmetric density by scaling its two
halves differently with a parameter $\xi > 0$:
\begin{equation}
  f_{\xi}(z; \nu)
  \;=\;
  \frac{2}{\xi + 1/\xi}\;
  f_{\mathrm{GED}}\!\left( \xi^{-\operatorname{sign}(z)}\, z ;\ \nu \right).
  \label{eq:fs}
\end{equation}
$\xi > 1$ stretches the right tail (positive skew), $\xi < 1$ the left;
$\xi = 1$ recovers symmetry. Skewing shifts the mean and changes the
variance, so \eqref{eq:fs} is not yet a location-scale-friendly object.

\paragraph{Step 3: re-standardization.}
With $\lambda = \lambda_\nu$, the absolute moments of the GED give
\begin{equation}
  m_1 = \frac{\Gamma(2/\nu)\,\lambda\, 2^{1/\nu}}{\Gamma(1/\nu)},
  \qquad
  m_2 = \frac{\Gamma(3/\nu)\,\lambda^2\, 2^{2/\nu}}{\Gamma(1/\nu)},
  \label{eq:sged-moments}
\end{equation}
from which the mean and variance of the skewed density \eqref{eq:fs} are
\begin{equation}
  \mu_\xi = \left(\xi - \tfrac{1}{\xi}\right) m_1,
  \qquad
  \sigma_\xi^2
  = \left(\xi^2 + \tfrac{1}{\xi^2}\right)\left(m_2 - m_1^2\right)
    + 2 m_1^2 - m_2 .
  \label{eq:sged-muxi}
\end{equation}
The standardized SGED --- mean $0$, variance $1$, the form used throughout
the thesis and the \code{fGarch} R package --- is then
\begin{equation}
  f^{\ast}_{\mathrm{SGED}}(z; \nu, \xi)
  \;=\;
  \sigma_\xi\, f_{\xi}\!\left( \sigma_\xi z + \mu_\xi ;\ \nu \right),
  \label{eq:sged-std}
\end{equation}
and the four-parameter jump law with location \code{jump\_loc} $= \ell$
and scale \code{jump\_scale} $= s$ (which \emph{is} the standard
deviation) is $f_J(x) = f^{\ast}_{\mathrm{SGED}}((x - \ell)/s;\, \nu, \xi)/s$.

\begin{lstlisting}[caption={The three-step SGED construction in \code{distributions/sged.py}.}]
def _ged_lambda(nu):
    return np.sqrt(2 ** (-2 / nu) * gamma(1 / nu) / gamma(3 / nu))

def _ged_pdf(z, nu):                     # eq. (GED): standardized base
    lam = _ged_lambda(nu)
    return (nu * np.exp(-0.5 * np.abs(z / lam) ** nu)
            / (lam * 2 ** (1 + 1 / nu) * gamma(1 / nu)))

def _fernandez_steel_pdf(z, nu, xi):     # eq. (FS): two-piece skewing
    sign = np.sign(z)
    return (2.0 / (xi + 1.0 / xi)) * _ged_pdf(np.power(xi, -sign) * z, nu)

def _skew_moments(nu, xi):               # eqs. (m1,m2) and (mu_xi, sigma_xi)
    lam = _ged_lambda(nu)
    m1 = gamma(2 / nu) * lam * 2 ** (1 / nu) / gamma(1 / nu)
    m2 = gamma(3 / nu) * lam**2 * 2 ** (2 / nu) / gamma(1 / nu)
    mu_xi = (xi - 1 / xi) * m1
    sigma_xi2 = (xi**2 + 1 / xi**2) * (m2 - m1**2) + 2 * m1**2 - m2
    return mu_xi, np.sqrt(sigma_xi2)

def _standardized_sged_pdf(z, nu, xi):   # eq. (SGED): mean 0, var 1
    mu_xi, sigma_xi = _skew_moments(nu, xi)
    return sigma_xi * _fernandez_steel_pdf(sigma_xi * z + mu_xi, nu, xi)
\end{lstlisting}

No closed form is known for $\phi * f_{\mathrm{SGED}}$, so this family
always goes through the FFT convolution of Section~\ref{sec:fft}. Its
special cases ($\nu{=}2, \xi{=}1 \Rightarrow$ normal; $\nu{=}1, \xi{=}1
\Rightarrow$ Laplace) are verified against closed forms in the test suite
--- a pattern used throughout the library: every numerical fallback is
tested where an exact answer exists.

\subsection{Kou's double-exponential jumps}
\label{sec:kou}

\citet{Kou2002} models jumps as a mixture of two one-sided exponentials:
with probability $q$ (\code{jump\_prob\_up}) the jump is positive,
$J \sim \mathrm{Exp}(\eta_u)$ with mean $\eta_u$
(\code{jump\_scale\_up}); otherwise it is negative with magnitude
$\sim \mathrm{Exp}(\eta_d)$ (\code{jump\_scale\_down}):
\begin{equation}
  f_J(x)
  \;=\;
  q\, \frac{1}{\eta_u} e^{-x/\eta_u}\, \Ind{x > 0}
  \;+\;
  (1 - q)\, \frac{1}{\eta_d} e^{x/\eta_d}\, \Ind{x < 0}.
  \label{eq:kou-pdf}
\end{equation}
Asymmetry here comes from $q \ne \tfrac12$ or $\eta_u \ne \eta_d$, not
from a shape parameter.

The Gaussian convolution again has a closed form, through the
\emph{exponentially modified Gaussian} (ex-Gaussian): if
$Y \sim \mathrm{Exp}(\eta)$ and $D \sim \mathcal N(m_d, \sigma_d^2)$,
the density of $D + Y$ is
\begin{equation}
  f_{D+Y}(x)
  =
  \frac{1}{2\eta}
  \exp\!\left( \frac{\sigma_d^2}{2\eta^2} - \frac{x - m_d}{\eta} \right)
  \operatorname{erfc}\!\left(
    \frac{1}{\sqrt2}\left( \frac{\sigma_d}{\eta} - \frac{x - m_d}{\sigma_d} \right)
  \right),
  \label{eq:exgauss}
\end{equation}
and the downward branch follows by the reflection $x \mapsto -x$. The
convolution of \eqref{eq:kou-pdf} with the diffusion is the
$q$-weighted mixture of one direct and one reflected ex-Gaussian.

\begin{decision}[Reuse \code{scipy.stats.exponnorm}, do not re-derive]
\label{dec:exponnorm}
Formula \eqref{eq:exgauss} is numerically treacherous: the exponential
factor overflows exactly where the $\operatorname{erfc}$ factor
underflows, and the stable evaluation requires the scaled complementary
error function $\operatorname{erfcx}$. SciPy's \code{exponnorm}
distribution already implements exactly this (parameterized by
$K = \eta / \sigma_d$), so the library delegates to it rather than
re-implementing the stability tricks --- and verifies the result against
direct numerical convolution in the tests.
\end{decision}

\begin{lstlisting}[caption={Kou convolution as a mixture of ex-Gaussians in \code{distributions/kou.py}.}]
def diffusion_convolved_pdf(self, x, params, diffusion_mean, diffusion_std):
    p = params["jump_prob_up"]
    up = p * exponnorm.pdf(
        x, K=params["jump_scale_up"] / diffusion_std,
        loc=diffusion_mean, scale=diffusion_std)
    down = (1 - p) * exponnorm.pdf(
        -x, K=params["jump_scale_down"] / diffusion_std,
        loc=-diffusion_mean, scale=diffusion_std)
    return up + down
\end{lstlisting}

Note the downward branch: reflecting the density requires negating
\emph{both} the evaluation points and the location of the Gaussian
component, which is why the code passes \code{-x} and
\code{loc=-diffusion\_mean}.

\subsection{Student-\texorpdfstring{$t$}{t} jumps}
\label{sec:student}

The heavy-tailed symmetric alternative: a location-scale Student-$t$ with
degrees of freedom $\nu_t$ (\code{jump\_df}), location \code{jump\_loc}
and scale \code{jump\_scale}.

\begin{decision}[Standardize the scale to the standard deviation]
\label{dec:t-scale}
SciPy's raw \code{t} scale $s_{\mathrm{scipy}}$ relates to the standard
deviation through the degrees of freedom:
$\operatorname{sd} = s_{\mathrm{scipy}} \sqrt{\nu_t / (\nu_t - 2)}$ for
$\nu_t > 2$. If the raw scale were exposed as the parameter, the fitted
``scale'' would not be comparable across candidate jump laws, and
$\nu_t$ would entangle spread with tail weight. The library therefore
re-parameterizes so that \code{jump\_scale} \emph{is} the standard
deviation,
\begin{equation}
  s_{\mathrm{scipy}} \;=\; s \sqrt{\frac{\nu_t - 2}{\nu_t}},
  \label{eq:t-scale}
\end{equation}
the same standardized-$t$ convention used for GARCH-$t$ innovations by
\citet{Bollerslev1987}. The lower bound of \code{jump\_df} is $2.05$,
keeping the variance finite so that \eqref{eq:t-scale} is well defined.
\end{decision}

No closed form exists for $\phi * t_{\nu_t}$, so like the SGED this family
relies on the FFT fallback --- which is the subject of the next section.

\section{Approximating the convolution by FFT}
\label{sec:fft}

When the jump family has no closed-form Gaussian convolution (SGED,
Student-$t$, and any user-added law that implements only \code{pdf}), the
mixture density \eqref{eq:mixture} needs the convolution
\begin{equation}
  (\phi_d * f_J)(x)
  \;=\;
  \int_{-\infty}^{\infty} \phi_d(t)\, f_J(x - t) \dd t,
  \qquad
  \phi_d \equiv \phi(\cdot\,;\ m_d, \sigma_d^2),
  \label{eq:conv-integral}
\end{equation}
evaluated at every observed increment, thousands of times per likelihood
maximization. The library approximates it with the discrete Fourier
transform, porting the scheme of \citet[ch.~4]{Ospina2009} (the
\code{f.rec} routine, originally in R) to NumPy. This section derives the
scheme and each of its tuning constants.

\subsection{From integral to circular convolution}

Discretize \eqref{eq:conv-integral} by the rectangle rule on a uniform
grid of spacing $h$: writing $x_k = k h$ for grid points,
\begin{equation}
  (\phi_d * f_J)(x_n)
  \;\approx\;
  h \sum_{k} \phi_d(x_k)\, f_J(x_n - x_k),
  \label{eq:riemann}
\end{equation}
which is $h$ times the \emph{discrete} (linear) convolution of the two
sampled densities. The FFT computes circular, not linear, convolution:
for two length-$N$ vectors $a$ and $b$,
\begin{equation}
  \operatorname{IDFT}\bigl( \operatorname{DFT}(a) \cdot \operatorname{DFT}(b) \bigr)_n
  \;=\;
  \sum_{k=0}^{N-1} a_k\, b_{(n - k) \bmod N},
  \label{eq:circular}
\end{equation}
at cost $O(N \log N)$ instead of $O(N^2)$. Circularity means the result
wraps around: mass that the linear convolution would place beyond the grid
ends re-enters at the opposite end. The scheme therefore only works if the
grid is wide enough that both densities are numerically negligible near
its ends --- then the wrapped-around contamination is negligible too. This
is the requirement that drives the grid-size decisions below.

\begin{remark}[NumPy vs.\ R normalization]
\label{rem:ifft}
Equation \eqref{eq:circular} holds for NumPy's convention, where
\code{np.fft.ifft} divides by $N$. R's \code{fft(..., inverse = TRUE)}
does \emph{not} divide by $N$, so the thesis' original R code carried an
explicit $1/N$ that must \emph{not} be ported. The code comments this
trap explicitly --- an example of a design decision that exists purely to
protect future maintainers.
\end{remark}

\subsection{The grid}
\label{sec:fft-grid}

The library samples both densities at $N = 2m$ points with
\emph{half-integer offsets},
\begin{equation}
  x_k \;=\; \left(k + \tfrac12\right) h,
  \qquad
  k = -m, -m+1, \dots, m-1,
  \qquad
  m = 2^{j},
  \label{eq:grid}
\end{equation}
i.e.\ a symmetric grid from $-(m - \tfrac12)h$ to $(m - \tfrac12)h$ that
straddles zero without containing it.

\begin{decision}[Half-integer offsets]
\label{dec:half-integer}
Two properties motivate the offset grid. First, symmetry: the grid covers
$[-mh, mh]$ evenly with no duplicated or privileged center point, so
symmetric densities are sampled symmetrically. Second, it dodges the
center singularity: some candidate jump densities are unbounded at their
mode (the Laplace corner at $0$, or the SGED with $\nu < 1$), and
evaluating exactly at $0$ would put one enormous (or infinite) sample into
the discrete convolution. With half-integer offsets every evaluation point
is at distance $\ge h/2$ from zero. The same grid construction is reused
by the generic sampler of Section~\ref{sec:contract}.
\end{decision}

\begin{decision}[Resolution first, coverage as a binding guard]
\label{dec:h-heuristic}
The step $h$ must resolve the \emph{narrower} of the two densities being
convolved --- oversampling the wide one costs nothing, undersampling the
narrow one destroys the approximation. Hence the base heuristic
\[
  h \;=\; \frac{\min(\sigma_d,\ s_J)}{200},
\]
where $s_J$ is the jump law's \code{characteristic\_scale}: two hundred
points per standard deviation of the narrowest component, with $j = 15$
giving $N = 2^{16} = 65{,}536$ grid points. Both constants are ported from
\citet{Ospina2009} and exposed as arguments (\code{j}, \code{h}).

Resolution alone, however, does not guarantee \emph{coverage}. The
half-width bought by the base heuristic is $mh \ge 163$ standard
deviations of the \emph{narrower} density --- which says nothing about
the \emph{wider} one. An optimizer legitimately probes candidates where
the jump scale dwarfs the diffusion scale (ratio $160{:}1$ leaves the
jump density with essentially zero mass on a grid sized by $\sigma_d$),
or where the jump location $\ell_J$ (the law's
\code{characteristic\_location}, Section~\ref{sec:contract}) sits far
from the origin. The step is therefore \emph{coarsened} whenever the span
could not contain the mass of both densities:
\begin{equation}
  m h \;<\; \underbrace{|m_d| + |\ell_J| + 10\,(\sigma_d + s_J)}_{\text{half-span needed}}
  \quad\Longrightarrow\quad
  h \;=\; \frac{|m_d| + |\ell_J| + 10\,(\sigma_d + s_J)}{m}.
  \label{eq:coarsen}
\end{equation}
A slightly coarser grid that preserves total mass beats a fine grid that
silently truncates it. In the default daily-data regime (scale ratios up
to $\approx 18$) the guard never binds, so ordinary fits are unaffected;
it exists for the extreme candidates that global optimizers and boundary
studies visit.
\end{decision}

\subsection{Re-centering and evaluation}

One subtlety remains. The circular convolution \eqref{eq:circular} indexes
vectors by array position $0, \dots, N-1$, while our samples live on the
physical grid \eqref{eq:grid}. If entry $i$ of each input vector holds the
density at $x_i = (i - m + \tfrac12) h$, then output entry $n$ of
\eqref{eq:circular} accumulates products of samples whose physical
locations add to
\[
  x_i + x_k \;=\; (i + k - 2m + 1)\, h ,
\]
with $i + k \equiv n \pmod N$. Matching this back to the target grid
\eqref{eq:grid} shows the output is displaced by half the array length:
entry $n$ of the raw FFT result corresponds to physical point
$x_{n - m}$ (modulo $N$), not $x_n$. Swapping the two halves of the output
array --- a circular shift by $m$, NumPy's \code{fftshift} --- restores
alignment, after which entry $n$ approximates
$(\phi_d * f_J)\bigl( x_n \bigr)$ on the same symmetric grid the inputs
were sampled on.

Finally, the convolved density is only known on grid points, and the
likelihood needs it at the observed increments: the code linearly
interpolates and clips tiny negative values produced by floating-point
cancellation to zero.

\begin{decision}[Out-of-grid points evaluate to zero --- individually]
\label{dec:out-of-grid}
Once the coverage guard \eqref{eq:coarsen} ensures the grid contains the
densities' mass, the true convolved density beyond the grid ends is
genuinely negligible --- so any observed increment falling there is
assigned density $0$ (\code{np.interp} with \code{left=0.0, right=0.0}),
which the likelihood's underflow floor (Section~\ref{sec:model}) then
turns into a strong but finite penalty on that parameter vector.

The qualifier \emph{individually} earns its place in the title: an
earlier version of this routine returned an all-zeros vector as soon as
\emph{any single} requested point fell outside the grid. The 2026 audit
of the library showed the consequence: one extreme observation (e.g.\ a
log-return of $0.11$ in the S\&P~500 sample) silently zeroed the whole
convolution term for \emph{every} candidate whose jump scale was small
enough that the grid, sized by Decision~\ref{dec:h-heuristic}, did not
reach that observation --- producing an artificially flat, floored
likelihood over the entire small-scale region. That cliff biased
estimation against small jump scales, broke the smoothness L-BFGS-B
relies on, and would have contaminated precisely the
$p \to 0$ / scale $\to 0$ boundary regions that the inference and
testing machinery of Sections~\ref{sec:inference}--\ref{sec:jumptest}
cares most about. The lesson generalizes: a numerical guard must fail
\emph{pointwise}, never collectively.
\end{decision}

\begin{lstlisting}[caption={The full FFT convolution in \code{distributions/base.py}.}]
def fft_convolved_pdf(self, x, params, diffusion_mean, diffusion_std,
                      j=15, h=None):
    x = np.asarray(x, dtype=float)
    jump_std = self.characteristic_scale(params)
    if jump_std is None or jump_std <= 0:
        return np.zeros_like(x)
    jump_loc = self.characteristic_location(params)

    m = 2**j
    if h is not None:
        step = h
    else:
        # Step chosen for resolution, then coarsened if the span could
        # not contain the mass of both densities (plus tail margins).
        step = min(diffusion_std, jump_std) / 200.0
        half_span_needed = (
            abs(diffusion_mean) + abs(jump_loc) + 10.0 * (diffusion_std + jump_std)
        )
        if m * step < half_span_needed:
            step = half_span_needed / m

    k = np.arange(-m + 0.5, m, 1.0)   # length 2*m, half-integer offsets
    x_grid = k * step

    f_diffusion = norm.pdf(x_grid, loc=diffusion_mean, scale=diffusion_std)
    f_jump = self.pdf(x_grid, params)
    if not (np.all(np.isfinite(f_diffusion)) and np.all(np.isfinite(f_jump))):
        return np.zeros_like(x)

    # numpy's ifft already divides by n, unlike R's fft(..., inverse=TRUE)
    conv = np.real(np.fft.ifft(np.fft.fft(f_diffusion) * np.fft.fft(f_jump)))
    conv *= step                                  # rectangle-rule weight
    conv = np.concatenate([conv[m:], conv[:m]])   # re-center around x_grid
    conv = np.maximum(conv, 0.0)

    return np.interp(x, x_grid, conv, left=0.0, right=0.0)
\end{lstlisting}

Line by line against the derivation: the grid is \eqref{eq:grid}; the
step selection implements Decision~\ref{dec:h-heuristic} with the
coverage guard \eqref{eq:coarsen}; the product of FFTs and single inverse
FFT is \eqref{eq:circular} (with the normalization of
Remark~\ref{rem:ifft}); the multiplication by \code{step} is the
rectangle-rule weight of \eqref{eq:riemann}; the \code{concatenate} is the
circular shift by $m$; and the final \code{interp}, with its pointwise
zeros outside the grid (Decision~\ref{dec:out-of-grid}), evaluates the
result at the observed increments.

\begin{remark}[Accuracy checks]
The scheme is validated in the test suite by running it on families whose
convolution \emph{is} known --- normal and skew-normal --- and comparing
against \eqref{eq:normal-conv} and \eqref{eq:sn-conv}, as well as by
checking the SGED special cases of Section~\ref{sec:sged}. This gives
end-to-end assurance without ever trusting the FFT code on faith for the
families where no closed form exists.
\end{remark}

\section{Maximum likelihood estimation}
\label{sec:estimation}

With the transition density in hand, estimation is the maximization of
\eqref{eq:loglik} over
$\thetavec = (\mu, \sigma, p, \theta_J)$. The estimator class,
\code{JumpDiffusionEstimator}, wraps this in a numerically defensive
objective and offers two optimizers whose contrast is one of the library's
central themes: a fast local method that needs a good start, and a global
method that needs none.

\subsection{The objective}

Optimizers minimize, so the objective is the \emph{negative}
log-likelihood, with two guard layers wrapped around \eqref{eq:loglik}:

\begin{lstlisting}[caption={The objective in \code{estimation/maximum\_likelihood.py}.}]
def log_likelihood(self, params: np.ndarray) -> float:
    sigma, jump_prob = params[1], params[2]

    if sigma <= 0:
        return np.inf
    if jump_prob < 0 or jump_prob > 1:
        return np.inf

    self._model.update_parameters(**dict(zip(self._param_names, params)))
    value = self._model.log_likelihood(self.increments, self.dt)
    # Penalize non-numeric likelihood values so that population-based
    # optimizers simply discard the offending candidates -- same device
    # as in Ospina Arango (2009).
    if not np.isfinite(value):
        return np.inf
    return -value
\end{lstlisting}

\begin{decision}[Map every pathology to $+\infty$]
\label{dec:penalize}
Invalid parameters ($\sigma \le 0$, $p \notin [0,1]$) and non-finite
likelihood evaluations (numerical accidents deep in the FFT or in extreme
tails) all return $+\infty$. For a population-based optimizer this is
exactly the right poison: an infinite objective makes the candidate lose
every selection comparison and disappear from the population without
corrupting the ranking of the finite candidates. The same device ---
mapping \code{NaN} evaluations to a large positive constant --- was used
in \citet{Ospina2009}; $+\infty$ is its cleanest float representation.
\end{decision}

\subsection{Starting values from moments}
\label{sec:initial-guess}

Gradient methods need a starting point. The estimator builds one from the
sample moments of the increments: since
$\E[\Delta X] = \mu \Delta t + p\,\E[J]$ and
$\operatorname{Var}[\Delta X] = \sigma^2 \Delta t + p \operatorname{Var}[J]
+ p(1-p)\E[J]^2$, ignoring the (typically small) jump contributions gives
the first-pass guesses
\begin{equation}
  \hat\mu_0 = \frac{\overline{\Delta x}}{\Delta t},
  \qquad
  \hat\sigma_0 = \frac{s_{\Delta x}}{\sqrt{\Delta t}},
  \qquad
  \hat p_0 = 0.1,
  \label{eq:initial-guess}
\end{equation}
and each jump family contributes its own heuristic through its
\code{initial\_guess} method --- e.g.\ the skew-normal
starts \code{jump\_scale} at half the increment standard deviation and
signs \code{jump\_skew} by the sample skewness. These are deliberately
crude: their job is to land in the right order of magnitude, not to be
accurate. When they are not enough --- and Section~\ref{sec:de} shows this
happens in practice --- the answer is a global optimizer, not a cleverer
guess.

\subsection{Local optimization: L-BFGS-B}

The default method hands the objective to
\code{scipy.optimize.minimize} with \code{method="L-BFGS-B"}, box
constraints from the model ($\sigma > 0$, $0 < p < 1$, plus each jump
parameter's bounds from Section~\ref{sec:contract}), and tolerances
\code{maxiter=1000}, \code{ftol=1e-9}. On well-behaved likelihoods (normal
or skew-normal jumps, moderate $p$) it converges in milliseconds. Its
failure modes on this problem are well documented in
\citet{Ospina2009}: with heavy-tailed or strongly skewed jump laws the
mixture likelihood develops multiple local optima and long, curved ridges,
and a quasi-Newton iterate either stalls on a box boundary or polishes
whichever local optimum captures it.

\subsection{Global optimization: differential evolution}
\label{sec:de}

Differential evolution \citep{StornPrice1997} maintains a population of
$\mathrm{NP}$ candidate vectors and evolves it by mutation, crossover and
greedy selection. One generation of the \emph{rand/1/bin} strategy --- the
one used here --- transforms each population member $x_i$ as follows:
\begin{align}
  \text{mutation:}\quad
    & v_i = x_{r_1} + F\,(x_{r_2} - x_{r_3}),
    \qquad r_1 \ne r_2 \ne r_3 \ne i \ \text{random},
    \label{eq:de-mutation}\\[2pt]
  \text{crossover:}\quad
    & u_{i,d} =
      \begin{cases}
        v_{i,d} & \text{if } U_d \le \mathrm{CR} \ \text{or}\ d = d_{\mathrm{rand}},\\
        x_{i,d} & \text{otherwise},
      \end{cases}
    \label{eq:de-crossover}\\[2pt]
  \text{selection:}\quad
    & x_i \leftarrow u_i \ \text{iff}\ f(u_i) \le f(x_i).
    \label{eq:de-selection}
\end{align}
Mutation \eqref{eq:de-mutation} perturbs a random member by a scaled
difference of two others, so step sizes adapt to the population's own
spread in each direction --- large early on, shrinking as the population
contracts onto a basin. Binomial crossover \eqref{eq:de-crossover} mixes
mutant and parent coordinate-wise. Selection \eqref{eq:de-selection} is
greedy but elitist: the best-ever point can never be lost. No gradients
appear anywhere, so multimodality, ridges, and the occasional
$+\infty$ from Decision~\ref{dec:penalize} are all handled gracefully.

\begin{decision}[Port the thesis' \code{DEoptim} configuration]
\label{dec:de-config}
The applied finding of \citet{Ospina2009} --- refined in
\citet{Ardia2011} --- is that DE with R's \code{DEoptim} defaults
(strategy rand/1, $\mathrm{NP} = 70$, 400 generations) recovers the SGED
jump-diffusion optimum even from very wide bounds, where L-BFGS-B and
simulated annealing both fail. The library ports this configuration to
\code{scipy.optimize.differential\_evolution}:
\begin{itemize}
  \item \code{strategy="rand1bin"} is exactly rand/1 with binomial
        crossover;
  \item SciPy sizes the population as $\mathrm{popsize} \times k$ where
        $k$ is the number of parameters, so \code{popsize=10} reproduces
        $\mathrm{NP} = 70$ for the 7-parameter SGED model and scales
        proportionally for other jump families;
  \item \code{maxiter=400} matches the generation cap, and SciPy's
        convergence tolerance (\code{tol=0.01}, stop when the population's
        objective spread falls below 1\% of its mean) implements the
        early-stopping criterion the thesis recommends --- in practice
        runs stop well before 400 generations.
\end{itemize}
\end{decision}

\begin{decision}[Finite, data-driven search box]
\label{dec:finite-bounds}
DE samples its initial population uniformly over the bounds, so the
half-open boxes used by L-BFGS-B ($\mu$ unbounded, $\sigma < \infty$)
must be replaced by finite ones. The library builds them from the data:
$\mu$ ranges over $\hat\mu_0 \pm \max(1,\ 10|\hat\mu_0|,\ 10\hat\sigma_0)$,
$\sigma$ over $(10^{-6},\ 10\,\hat\sigma_0)$, and any unbounded jump
parameter is capped at $\pm 20\, s_{\Delta x}$ --- a single jump twenty
times the typical increment is a generous ceiling, especially since
$s_{\Delta x}$ is itself inflated by whatever jumps the data contains.
Wide boxes are deliberate: the thesis' point is precisely that DE does
not need tight prior knowledge, so the defaults encode ignorance, not
information. Bounds the user supplies, and finite bounds from the jump
family, are kept as-is.
\end{decision}

\begin{lstlisting}[caption={DE configuration in \code{estimation/maximum\_likelihood.py} (excerpt).}]
if bounds is None:
    bounds = self._finite_bounds()
lows, highs = zip(*bounds)
de_bounds = Bounds(np.asarray(lows, float), np.asarray(highs, float))

de_kwargs = {
    "strategy": "rand1bin",   # DEoptim's rand/1 with binomial crossover
    "maxiter": 400,           # thesis generation cap
    "popsize": 10,            # 10 * 7 params = NP 70 for the SGED model
}
de_kwargs.update(kwargs)      # user overrides, e.g. seed=, workers=
return differential_evolution(func=self.log_likelihood,
                              bounds=de_bounds, **de_kwargs)
\end{lstlisting}

The cost asymmetry is worth stating plainly: L-BFGS-B typically needs
tens-to-hundreds of likelihood evaluations, DE thousands --- seconds
instead of milliseconds with closed-form convolutions, and proportionally
more through the FFT path. The recommended workflow is therefore to fit
with DE (optionally passing \code{seed=} for reproducibility and
\code{workers=-1} to parallelize the population evaluation) and reserve
L-BFGS-B for quick exploratory fits and for the repeated re-fits inside
bootstrap loops, where the data are simulated near a known optimum.

\subsection{Reported quantities}

\code{estimate()} returns the fitted parameters, the maximized
log-likelihood $\loglik(\hat\thetavec)$, and the information criteria
\begin{equation}
  \mathrm{AIC} = 2k - 2\loglik(\hat\thetavec),
  \qquad
  \mathrm{BIC} = k \log n - 2\loglik(\hat\thetavec),
  \label{eq:aic-bic}
\end{equation}
with $k = \dim\thetavec$ and $n$ the number of increments --- the
ingredients for the model comparison of Section~\ref{sec:comparison} ---
plus the raw optimizer diagnostics and a convergence flag that downstream
bootstrap loops use to discard failed re-fits.

\section{Uncertainty quantification: three routes}
\label{sec:inference}

A point estimate without standard errors is half an answer. The library
implements three inference routes with deliberately different assumptions
and costs, summarized in Table~\ref{tab:inference-routes}. Offering all
three is itself a design decision: comparing their finite-sample coverage
across jump families is one of the research questions this library was
built to answer (Section~\ref{sec:closing}).

\begin{table}[ht]
\centering
\small
\begin{tabular}{@{}p{5.6cm}p{4.5cm}l@{}}
\toprule
Route (method) & Assumes & Cost \\
\midrule
Wald\newline {\footnotesize\code{estimate\_wald\_standard\_errors}} &
  asymptotic normality, regular MLE & one Hessian \\
\addlinespace
Profile\newline {\footnotesize\code{estimate\_standard\_errors}} &
  Wilks' $\chi^2$ asymptotics & re-optimizations \\
\addlinespace
Bootstrap\newline {\footnotesize\code{estimate\_bootstrap\_standard\_errors}} &
  correct model specification & $B$ full re-fits \\
\bottomrule
\end{tabular}
\caption{The three inference routes. Cost grows downward; robustness to
boundary and non-quadratic likelihoods grows the same way.}
\label{tab:inference-routes}
\end{table}

\subsection{Wald intervals from the observed information}
\label{sec:wald}

Classical first-order asymptotics: under regularity conditions,
$\hat\thetavec \approx \mathcal N\bigl(\thetavec_0,\ I(\hat\thetavec)^{-1}\bigr)$
where
\begin{equation}
  I(\hat\thetavec) \;=\; -\nabla^2 \loglik(\hat\thetavec)
  \label{eq:observed-info}
\end{equation}
is the \emph{observed} Fisher information --- the Hessian of the negative
log-likelihood at the estimate. Standard errors are the square roots of
the diagonal of $I^{-1}$, and intervals are
$\hat\theta_r \pm z_{1-\alpha/2}\,\widehat{\mathrm{SE}}_r$.

\begin{decision}[Observed, not expected, information]
Following \citet{EfronHinkley1978}, the library uses the observed
information rather than the expected Fisher information: for this mixture
likelihood the expectation has no closed form and would itself require
numerical integration, and Efron and Hinkley argue the observed
information is the more relevant conditioning quantity anyway.
\end{decision}

Since the code's objective is already the negative log-likelihood, its
Hessian is exactly \eqref{eq:observed-info}, built by central finite
differences: on the diagonal
\begin{equation}
  \frac{\partial^2 f}{\partial \theta_r^2}
  \;\approx\;
  \frac{f(\hat\thetavec + h_r e_r) - 2 f(\hat\thetavec)
        + f(\hat\thetavec - h_r e_r)}{h_r^2},
  \label{eq:fd-diag}
\end{equation}
and off the diagonal the four-point mixed formula
\begin{equation}
  \frac{\partial^2 f}{\partial \theta_r \partial \theta_s}
  \;\approx\;
  \frac{f_{++} - f_{+-} - f_{-+} + f_{--}}{4 h_r h_s},
  \label{eq:fd-mixed}
\end{equation}
where $f_{\pm\pm} = f(\hat\thetavec \pm h_r e_r \pm h_s e_s)$.

\begin{decision}[Bound-aware step sizes]
\label{dec:fd-steps}
The step for parameter $r$ is $h_r = 10^{-4} \max(|\hat\theta_r|, 10^{-2})$
--- relative, with an absolute floor so parameters estimated near zero
(drift, jump location) still get a usable step. Crucially, each step is
then shrunk to at most a quarter of the distance to the nearest bound:
the objective returns $+\infty$ outside the bounds
(Decision~\ref{dec:penalize}), and a single infinite evaluation would
destroy the difference quotients \eqref{eq:fd-diag}--\eqref{eq:fd-mixed}.
Entries that still cannot be evaluated (a parameter pinned exactly to a
bound) are recorded as \code{NaN}, and the inversion is guarded: if the
Hessian is non-finite or singular, every Wald SE is reported as
\code{NaN} rather than silently wrong. This is not an edge case ---
$\hat p$ near $0$ or a jump scale collapsing toward its lower bound are
exactly the interesting regimes.
\end{decision}

\subsection{Profile likelihood intervals}
\label{sec:profile}

The Wald construction forces symmetric intervals and trusts a quadratic
approximation of the log-likelihood. The profile likelihood does neither.
For a scalar parameter of interest $\theta_r$, define
\begin{equation}
  \loglik_p(v)
  \;=\;
  \max_{\thetavec\,:\,\theta_r = v} \loglik(\thetavec),
  \label{eq:profile}
\end{equation}
the log-likelihood maximized over all \emph{other} parameters with
$\theta_r$ pinned at $v$. By Wilks' theorem \citep{Wilks1938}, under the
true value $2[\loglik(\hat\thetavec) - \loglik_p(\theta_{r,0})]
\rightsquigarrow \chi^2_1$, so
\begin{equation}
  \mathrm{CI}_{1-\alpha}
  \;=\;
  \Bigl\{ v \;:\; \loglik_p(v) \;\ge\;
  \underbrace{\loglik(\hat\thetavec) - \tfrac12 \chi^2_{1,1-\alpha}}_{\text{threshold}}
  \Bigr\}
  \label{eq:profile-ci}
\end{equation}
is an asymptotic $(1-\alpha)$ confidence set that follows the likelihood's
actual shape --- asymmetric where the likelihood is asymmetric, respecting
parameter bounds by construction.

Each profile point \eqref{eq:profile} is a constrained re-optimization
(L-BFGS-B over the remaining parameters, warm-started at the MLE), so the
practical question is \emph{where} to spend those evaluations. The
library's algorithm, per parameter:

\begin{enumerate}
  \item \textbf{Adaptive outward search.} Starting from the MLE with step
        $\max(0.01\,|\hat\theta_r|,\ 10^{-4})$, walk right, doubling the
        step each time, evaluating $\loglik_p$ at each point, until the
        profile drops below the threshold in \eqref{eq:profile-ci}, a
        parameter bound is hit (evaluate just inside it), or 15 steps are
        exhausted. Repeat leftward. Geometric doubling reaches a crossing
        $k$ standard errors away in $O(\log k)$ evaluations --- important
        because parameters like $p$ can have profiles an order of
        magnitude wider than their Wald SE suggests.
  \item \textbf{Infill.} If fewer than \code{n\_points} evaluations
        accumulated, bisect the largest remaining gap between consecutive
        evaluated points, repeatedly --- concentrating resolution where
        the profile is least known.
  \item \textbf{Interval endpoints by interpolation.} The CI endpoints are
        where $\loglik_p$ crosses the threshold: found by linear
        interpolation between the bracketing evaluations on each side. If
        the profile never drops below the threshold before a bound (e.g.\
        $p$ against $0$), the interval is truncated at the searched range
        --- reproducing the one-sided behavior a boundary parameter should
        have.
  \item \textbf{Standard error from local curvature.} Near the MLE,
        $\loglik_p(v) \approx \loglik(\hat\thetavec) -
        (v - \hat\theta_r)^2 / (2\,\mathrm{SE}^2)$: a parabola whose
        leading coefficient $a$ satisfies
        $\mathrm{SE} = \sqrt{-1/(2a)}$. The code fits that parabola by
        least squares through the (up to five) highest-likelihood profile
        points. If the fit is not concave ($a \ge 0$, a genuinely
        non-quadratic profile), the fallback derives the SE from the
        interval width instead, $\mathrm{SE} \approx
        (\mathrm{CI}_{\mathrm{high}} - \mathrm{CI}_{\mathrm{low}}) /
        (2 z_{1-\alpha/2})$ --- and only when the endpoints genuinely
        crossed the threshold rather than being truncated at the searched
        range.
\end{enumerate}

\begin{lstlisting}[caption={The profiled objective: pin $\theta_r$, optimize the rest (excerpt).}]
def optimize_profile(val):
    def neg_log_lik_profile(free_params):
        full_params = np.zeros(len(self._param_names))
        full_params[idx] = val          # parameter of interest, pinned
        free_idx = 0
        for j in range(len(self._param_names)):
            if j != idx:
                full_params[j] = free_params[free_idx]
                free_idx += 1
        return self.log_likelihood(full_params)

    res = minimize(fun=neg_log_lik_profile, x0=free_init,
                   method="L-BFGS-B", bounds=free_bounds,
                   options={"maxiter": 100, "ftol": 1e-6})
    return -res.fun if res.success and np.isfinite(res.fun) else -np.inf
\end{lstlisting}

The companion method \code{plot\_profiles()} draws each profile with its
threshold line --- the visual version of \eqref{eq:profile-ci}, and in
practice the fastest way to \emph{see} identifiability problems: a flat
profile is a parameter the data cannot pin down.

\subsection{Parametric bootstrap}
\label{sec:bootstrap}

The third route drops the asymptotics altogether and pays with
computation. Treat the fitted model as the truth; simulate $B$ datasets of
the same length $n$ and spacing $\Delta t$ from it; re-estimate
$\thetavec$ on each; and read the sampling distribution of the estimator
directly off the $B$ replicate estimates
$\hat\thetavec^{(1)}, \dots, \hat\thetavec^{(B)}$
\citep{DavisonHinkley1997}:
\begin{equation}
  \widehat{\mathrm{SE}}_r = \operatorname{sd}\bigl(\hat\theta_r^{(b)}\bigr),
  \qquad
  \mathrm{CI}_{1-\alpha} =
  \Bigl[ \hat\theta_r^{(\lceil B\alpha/2 \rceil)},\
         \hat\theta_r^{(\lceil B(1-\alpha/2) \rceil)} \Bigr]
  \quad \text{(percentile interval)}.
  \label{eq:bootstrap}
\end{equation}
No normality, no quadratic approximation, and boundary effects are
represented honestly: if $\hat p$ piles up near zero across replicates,
the percentile interval shows it.

Implementation decisions worth recording: replicate $b$ is simulated with
seed $\texttt{seed} + b$, making the entire bootstrap reproducible from
one integer while keeping replicates independent; re-fits that fail to
converge are discarded (their count is reported as
\code{n\_successful}), since averaging over failed optimizations would
contaminate \eqref{eq:bootstrap} with optimizer artifacts; the SE uses the
$\mathrm{ddof}=1$ sample standard deviation; and the raw replicate matrix
is stored in the results dictionary so users can compute BCa or basic
intervals, correlations between parameters, or replicate diagnostics
without re-running the loop.

\begin{lstlisting}[caption={The bootstrap loop (excerpt).}]
model = JumpDiffusionModel(jump_distribution=jump_dist, **fitted_params)
for b in range(n_bootstrap):
    rep_seed = None if seed is None else seed + b
    _, path, _ = model.simulate(T=horizon, n_steps=n_steps, x0=0.0,
                                seed=rep_seed)
    increments = np.diff(path)
    estimator = JumpDiffusionEstimator(increments, self.dt,
                                       jump_distribution=jump_dist)
    rep_result = estimator.estimate(**estimate_kwargs)
    if rep_result["convergence"]:
        replicates.append([rep_result["parameters"][name]
                           for name in self._param_names])
\end{lstlisting}

All three routes converge into one view: \code{summary()} returns a tidy
table with one row per parameter and one column block per computed route,
and \code{diagnostics()} prints it with convergence and data statistics
--- making the disagreement between routes (the interesting part) visible
at a glance.

\code{diagnostics()} also compares the sample mean and variance of the
increments against the ones the fitted model \emph{implies}. Under the
Bernoulli mixture \eqref{eq:bernoulli-increment} these are exact:
\begin{equation}
  \E[\Delta X] = \mu\,\Delta t + p\,\E[J],
  \qquad
  \operatorname{Var}[\Delta X]
  = \sigma^2 \Delta t + p\,\E[J^2] - \bigl(p\,\E[J]\bigr)^2,
  \label{eq:incr-moments}
\end{equation}
with $\E[J]$ and $\operatorname{Var}[J]$ supplied by the jump law's
generic numerical moments (Section~\ref{sec:contract}).

\begin{remark}[The jump term in the mean is not optional]
An earlier version of \code{diagnostics()} printed $\mu\,\Delta t$ alone
as the theoretical mean. With symmetric jumps that omission is invisible
($\E[J]=0$), but the 2026 audit measured its effect on a skew-normal fit:
a reported theoretical mean of $-0.000019$ against an empirical mean of
$0.011798$ --- suggesting a gross model--data mismatch --- when the
model's actual implied mean, $\mu\Delta t + p\,\E[J] = 0.011778$, agreed
with the data almost exactly. For \emph{asymmetric} jump laws, the whole
point of this library, the $p\,\E[J]$ term routinely dominates
$\mu\,\Delta t$; equation \eqref{eq:incr-moments} is now computed in
full.
\end{remark}

\section{Testing for the presence of jumps}
\label{sec:jumptest}

Before interpreting fitted jump parameters, one should ask whether the
data support jumps at all:
\begin{equation}
  H_0:\ p = 0
  \qquad \text{vs.} \qquad
  H_1:\ p > 0.
  \label{eq:jump-hypotheses}
\end{equation}
The natural statistic is the likelihood ratio. Under $H_0$ the model
collapses to pure diffusion, whose increments are i.i.d.\
$\mathcal N(\mu\Delta t, \sigma^2\Delta t)$ --- so the null MLE is
available in closed form (sample mean and MLE standard deviation of the
increments, no optimizer needed) and
\begin{equation}
  \mathrm{LR}
  \;=\;
  2\left[ \loglik(\hat\thetavec) - \loglik_0(\hat m, \hat s) \right],
  \qquad
  \loglik_0(\hat m, \hat s) = \sum_{i=1}^{n}
  \log \phi(\Delta x_i;\ \hat m,\ \hat s^2).
  \label{eq:lr}
\end{equation}

\subsection{Why the textbook calibration fails twice}

Were this a regular testing problem, $\mathrm{LR} \rightsquigarrow
\chi^2_{d}$ with $d$ the number of restricted parameters. Here two
pathologies invalidate that, \emph{simultaneously}:

\paragraph{(i) The null is on the boundary.} $p = 0$ is not an interior
point of $[0, 1]$. When a tested parameter sits on the boundary of the
parameter space, the LR statistic's limit is not $\chi^2$ but a
\emph{mixture} --- in the simplest case
$\tfrac12 \chi^2_0 + \tfrac12 \chi^2_1$, reflecting that with probability
one half the unconstrained estimate falls on the wrong side of the
boundary and the statistic degenerates to zero
\citep{Chernoff1954, SelfLiang1987}.

\paragraph{(ii) Nuisance parameters vanish under the null.} When $p = 0$,
the jump-size parameters $\theta_J$ multiply nothing: the likelihood
\eqref{eq:mixture} is flat in $\theta_J$, so they are \emph{unidentified}
under $H_0$. This is Davies' problem \citep{Davies1977, Davies1987}: the
LR statistic is then the supremum of a process indexed by the unidentified
parameters, its null distribution depends on the model and the data
design, and no universal table exists.

Each pathology alone has known corrections; their combination, for this
mixture with five different jump families, does not admit a practical
closed-form calibration.

\subsection{The parametric bootstrap test}

\begin{decision}[Simulate the null distribution instead of assuming one]
\label{dec:bootstrap-test}
The library sidesteps both pathologies with a Monte Carlo test
\citep{DavisonHinkley1997}: whatever the true null distribution of
$\mathrm{LR}$ is --- boundary mixture, Davies supremum, finite-sample
quirks and all --- it can be sampled from, because the null model is fully
specified once fitted. The procedure:
\begin{enumerate}
  \item Fit the null on the observed data: $\hat m, \hat s$ in closed
        form; compute the observed $\mathrm{LR}$ via \eqref{eq:lr}.
  \item For $b = 1, \dots, B$: simulate $n$ increments i.i.d.\
        $\mathcal N(\hat m, \hat s^2)$; fit the \emph{full} jump-diffusion
        model on them; compute $\mathrm{LR}^{(b)}$ the same way.
  \item Report the Monte Carlo p-value
        \begin{equation}
          \hat p_{\mathrm{MC}}
          \;=\;
          \frac{1 + \#\{ b :\ \mathrm{LR}^{(b)} \ge \mathrm{LR} \}}
               {B_{\mathrm{succ}} + 1},
          \label{eq:mc-pvalue}
        \end{equation}
        where $B_{\mathrm{succ}}$ counts the replicates whose full-model
        re-fit converged.
\end{enumerate}
The ``$+1$'' in numerator and denominator makes \eqref{eq:mc-pvalue} an
exact-level p-value: it treats the observed statistic as one more draw
from the null, which is precisely what it is when $H_0$ holds
\citep{DavisonHinkley1997}. The price is $B$ full model fits --- the same
cost profile as the bootstrap intervals of Section~\ref{sec:bootstrap}.
\end{decision}

Two implementation details guard the statistic itself. First, the full
model nests the null, so in exact arithmetic $\mathrm{LR} \ge 0$; in
practice the optimizer can stop a hair short of the null likelihood, so
the code clamps the statistic at zero --- on both the observed data and
every replicate, keeping the comparison in \eqref{eq:mc-pvalue}
consistent. Second, the null MLE uses the population ($\mathrm{ddof}=0$)
standard deviation: that is the Gaussian \emph{maximum likelihood}
estimate, and using the unbiased ($\mathrm{ddof}=1$) version would
slightly misstate $\loglik_0$ and bias the statistic.

\begin{lstlisting}[caption={The jump test in \code{estimation/maximum\_likelihood.py} (excerpt).}]
log_lik_full = results["log_likelihood"]
log_lik_null, mean0, std0 = self._pure_diffusion_log_likelihood(self.increments)
# The full model nests the null, so LR is non-negative in theory;
# clamp to guard against the optimizer stopping just short of it.
lr_observed = max(0.0, 2.0 * (log_lik_full - log_lik_null))

rng = np.random.default_rng(seed)
bootstrap_stats = []
for _ in range(n_bootstrap):
    increments_b = rng.normal(mean0, std0, size=self.n_obs)
    estimator = JumpDiffusionEstimator(increments_b, self.dt,
                                       jump_distribution=jump_dist)
    rep_result = estimator.estimate(**estimate_kwargs)
    if not rep_result["convergence"]:
        continue
    log_lik_full_b = rep_result["log_likelihood"]
    log_lik_null_b, _, _ = self._pure_diffusion_log_likelihood(increments_b)
    bootstrap_stats.append(max(0.0, 2.0 * (log_lik_full_b - log_lik_null_b)))

exceed = int(np.sum(np.asarray(bootstrap_stats) >= lr_observed))
p_value = (1.0 + exceed) / (len(bootstrap_stats) + 1.0)
\end{lstlisting}

\begin{remark}[What remains to be established]
The bootstrap test is valid by construction \emph{under the fitted null};
its finite-sample size across sample lengths and its power against
alternatives with small $p$ or jump scales close to the diffusion scale
are empirical questions --- precisely the ones the Monte Carlo study of
Section~\ref{sec:closing} is designed to answer, with the naive
$\chi^2$ and the $\tfrac12$--$\tfrac12$ mixture calibrations as
comparators.
\end{remark}

\section{Comparing jump distributions and validating the estimator}
\label{sec:comparison}

\subsection{Ranking candidate jump laws}

With five interchangeable jump families, the applied question becomes:
which one does \emph{this} dataset support?
\code{JumpDistributionComparison} fits each candidate to the same
increments and reports two complementary kinds of evidence.

\paragraph{Information criteria.} AIC and BIC from \eqref{eq:aic-bic}
penalize the maximized likelihood for parameter count --- relevant here
because the candidates differ in dimension (Merton normal has one jump
parameter, the SGED four). They answer a \emph{relative} question: which
candidate trades fit against complexity best.

\paragraph{Goodness of fit.} Information criteria cannot say whether even
the best candidate actually fits. For that, the comparison computes a
Kolmogorov--Smirnov distance between the observed increments and the
fitted model's increment distribution --- the latter having no closed-form
CDF, it is represented by a large simulated reference sample from the
fitted model, so the statistic is a two-sample KS distance whose
reference-side noise is made negligible by sample size.

\begin{decision}[A bootstrap p-value, because parameters were estimated]
\label{dec:ks-bootstrap}
The classical KS p-value assumes the reference distribution was fixed
\emph{a priori}. Here it was estimated from the very data being tested,
which makes the naive p-value strongly anti-conservative --- the fitted
model always looks too good, the same phenomenon that motivates the
Lilliefors correction for Gaussian fits. The library restores validity
with the same device as Section~\ref{sec:jumptest}: simulate
$B$ datasets from the fitted model, \emph{re-fit the model on each}
(re-fitting is what accounts for estimation error), recompute each
replicate's KS distance to its own re-fitted model, and report the Monte
Carlo p-value \eqref{eq:mc-pvalue}. A model is inadequate when the
observed distance is large relative to the distances the model produces
against itself.
\end{decision}

\begin{lstlisting}[caption={Bootstrap KS p-value in \code{validation/distribution\_comparison.py} (excerpt).}]
for b in range(n_bootstrap):
    boot_data = self._simulate_increments(fitted_model, n_obs, data_seed)
    estimator = JumpDiffusionEstimator(
        boot_data, self.dt, jump_distribution=jump_distribution)
    rep_result = estimator.estimate(**estimate_kwargs)
    if not rep_result["convergence"]:
        continue
    model_b = JumpDiffusionModel(jump_distribution=jump_distribution,
                                 **rep_result["parameters"])
    reference_b = self._simulate_increments(model_b, reference_size, ref_seed)
    ks_b = float(ks_2samp(boot_data, reference_b).statistic)
    if ks_b >= ks_observed:
        exceed += 1
    n_successful += 1

return (1.0 + exceed) / (n_successful + 1.0)
\end{lstlisting}

\code{compare()} assembles everything into a table ranked by AIC (with
BIC, KS statistic and bootstrap p-value alongside), and
\code{plot\_comparison()} overlays each fitted mixture density on the
increment histogram --- the visual complement to the numbers. Because the
KS machinery relies only on simulation through
\code{JumpDistribution.rvs}, it works unchanged for any user-added jump
family, honoring the extension contract of Section~\ref{sec:contract}.

\subsection{Monte Carlo validation of the estimator}

Finally, \code{ValidationExperiment} closes the loop on the estimator
itself: choose true parameters, simulate $R$ independent paths, estimate
on each, and summarize the sampling behavior of the estimator ---
bias, standard deviation, RMSE per parameter, along with convergence
rates (all summary statistics condition on the runs whose optimization
converged). Two reporting conventions are deliberate. Relative errors are
reported as \code{NaN} whenever the true parameter is (numerically) zero
--- $\mu = 0$, \code{jump\_skew} $= 0$ and \code{jump\_loc} $= 0$ are all
natural experimental settings, and division by zero would otherwise
inject $\pm\infty$ into every downstream summary; \code{NaN} says
``undefined'', which is the honest answer. And the share of runs landing
within 5\% of the truth is reported as an \emph{accuracy rate}
(\code{within\_5pct\_rate}), pointedly not called ``coverage'': coverage
belongs to the confidence intervals of Section~\ref{sec:inference}, and
conflating the two invites exactly the kind of misreading a validation
tool exists to prevent.

This simulate--estimate--compare pattern is how every method in this
paper was tested during development, packaged for users; it is also the
seed of the full factorial Monte Carlo study (coverage of the three
interval routes, size and power of the jump test, identifiability maps)
that constitutes the research program around the library.

\section{Closing remarks and research agenda}
\label{sec:closing}

This paper has walked through the full pipeline implemented by
\code{jump-diffusion-estimation}: an exactly-solvable SDE whose only
approximation is the Bernoulli treatment of the jump count
(Section~\ref{sec:model}); five jump families behind a density-only
extension contract, with closed-form Gaussian convolutions exploited
wherever the mathematics allows (Section~\ref{sec:distributions}) and an
FFT scheme, ported from \citet{Ospina2009}, wherever it does not
(Section~\ref{sec:fft}); maximum likelihood driven by either a local
quasi-Newton method or the differential evolution configuration whose
robustness was the thesis' applied finding (Section~\ref{sec:estimation});
three inference routes of increasing robustness and cost
(Section~\ref{sec:inference}); a parametric-bootstrap likelihood-ratio
test that sidesteps the boundary and Davies pathologies by simulating the
null distribution outright (Section~\ref{sec:jumptest}); and
estimation-error-aware goodness-of-fit comparison across jump families
(Section~\ref{sec:comparison}).

A recurring pattern deserves to be named, because it \emph{is} the
library's design philosophy:

\begin{itemize}
  \item \textbf{Closed forms first, generic numerics as fallback.} Every
        numerical fallback (FFT convolution, inverse-CDF sampling) is
        validated against the closed forms available in special cases.
  \item \textbf{Pathologies are mapped to explicit poison values} ---
        $+\infty$ objectives, \code{NaN} standard errors --- never
        silently absorbed.
  \item \textbf{When asymptotics are doubtful, simulate.} The bootstrap
        appears three times (intervals, jump test, KS p-value), each time
        replacing a questionable asymptotic calibration with an honest
        Monte Carlo one.
  \item \textbf{Reproducibility is an argument, not an accident:} every
        stochastic routine accepts a seed and derives per-replicate seeds
        deterministically from it.
\end{itemize}

\subsection*{Open questions}

The library makes a set of methods available; it does not yet certify
their finite-sample behavior. The companion research program, for which
this library is the experimental engine, targets exactly that gap:

\begin{enumerate}
  \item \textbf{Coverage.} Empirical coverage of the Wald, profile and
        bootstrap intervals across jump families, sample sizes, jump
        intensities and jump-to-diffusion scale ratios --- when does
        profiling recover the nominal coverage that Wald loses?
  \item \textbf{Jump test calibration.} Size and power of the bootstrap
        LR test of Section~\ref{sec:jumptest} against the naive
        $\chi^2$ and the boundary-corrected mixture calibrations.
  \item \textbf{Identifiability.} Mapping the region of
        $(\lambda\Delta t,\ \text{jump scale}/\sigma\sqrt{\Delta t})$
        where $p$, $\sigma$ and the jump scale can be jointly recovered,
        via the conditioning of the observed information and the sampling
        correlation of $(\hat p, \hat\sigma)$.
  \item \textbf{Bernoulli error.} Quantifying the bias of
        Decision~\ref{dec:bernoulli} against the exact Poisson mixture as
        $\lambda\Delta t$ grows.
  \item \textbf{Optimizer geography.} Where in parameter space multistart
        L-BFGS-B fails and differential evolution succeeds, extending
        \citet{Ardia2011} from a demonstration to a map.
\end{enumerate}

Each question already has its instrument in the library; what remains is
the systematic experiment. That study is deliberately kept out of this
document: here we fixed \emph{what the code computes and why}, so that the
experiments --- and any reader's own use of the library --- stand on
explicit mathematical ground.


\bibliographystyle{plainnat}
\bibliography{references}

\end{document}
